\documentclass[10pt,a4paper]{amsart}
\usepackage[margin=19mm]{geometry}
\usepackage{amsmath,amssymb,mathtools}
\usepackage[scr=rsfs,cal=euler]{mathalfa}
\usepackage{xfrac}
\usepackage[unicode,hidelinks,bookmarks=false]{hyperref}
\newcommand{\ie}{i.e.\ }
\newcommand{\cf}{cf.\ }
\newcommand{\resp}{respectively\ }
\newcommand{\Subsection}{\subsection}
\providecommand{\nobreakdash}{\nobreak\hbox{-}}
\setlength{\emergencystretch}{2em}
\multlinegap0pt
\usepackage{amssymb}
\newtheorem{lemma}{Lemma}
\newtheorem{theorem}{Theorem}
\newtheorem{corollary}{Corollary}
\newtheorem{proposition}{Proposition}
\newcommand{\mcS}{\mathcal{S}}
\newcommand{\mcT}{\mathcal{T}}
\newcommand{\mcV}{\mathcal{V}}
\newcommand{\pro}[1]{{#1}^p}
\title{Semilattice sums of algebras\\ and Mal'tsev products of varieties}
\author{Clifford Bergman, Tomasz Penza, Anna B. Romanowska}
\date{Source: arXiv:2603.29747v1 (CC BY 4.0)}
\begin{document}
\maketitle

\noindent\emph{Source and licence.} Semilattice sums of algebras\\ and Mal'tsev products of varieties. Version arXiv:2603.29747v1 (CC BY 4.0), distributed by arXiv under the Creative Commons Attribution 4.0 International licence, http://creativecommons.org/licenses/by/4.0/, as recorded in the arXiv licence metadata for this exact version and shown on the abstract page https://arxiv.org/abs/2603.29747v1. Full licence notice: \texttt{licenses/arxiv-2603.29747-CC-BY-4.0.txt}.\par\smallskip

\noindent\textit{Editorial scope.} Counted unit: the article's proposition P:quotientC (source0.tex lines 952--955). The article's complete original proof of it is delivered with it (source0.tex lines 956--1016, one \texttt{proof} environment of 61 lines). No mathematics is written, repaired or re-derived: the statement and the proof are the article's own lines, in the article's own notation, and no retained line is substituted or re-flowed.  Retained uncounted as the source-local context the proof needs: lines 673--680; lines 725--728; lines 940--948.  Nothing was omitted from any retained span, so the package carries no omission record.  No retained line contains a citation, so the wrapper carries no bibliography and no bibliography entry.  No retained line contains a figure environment, an image-inclusion command or drawing code, so no image is delivered and the package declares no asset dependency.  Editorial formatting only, in addition: an AMS \texttt{amsart} preamble that restates the article's own declarations and macros used by the retained lines, namely the environments \texttt{lemma}, \texttt{theorem}, \texttt{corollary}, \texttt{proposition} on separate counters, together with the standard packages those lines need.  The retained environments print in this document as Lemma 1, Theorem 1, Corollary 1, Proposition 1 in order of appearance, whereas the article prints the same items with the numbering of its own full document.  The complete native source of the article as distributed in the arXiv e-print for this exact version is bundled unchanged as \texttt{sources/197537/source0.tex}.
\begin{lemma}\label{L:equivlz}
Let $\theta$ be a congruence of a $\pro{\mcT_t}$-algebra $A$, and let $\psi = \theta \vee
\varrho$. Then a congruence class $(a_r/\theta)/(\psi/\theta)$  of $A/\theta$ satisfies
the identity $x \cdot y = x$ if and only if, for any $b_s \in A$
\begin{equation}\label{E:equivlz}
(a_r, b_s) \in \psi \implies (a_r b_s, a_r) \in \theta.
\end{equation}
\end{lemma}

\begin{theorem}\label{T:varPt}
The quasivariety $\mcT_{t} \circ \mcS$ and the variety $\pro{\mcT_t}$ coincide.
In particular, the Mal'tsev product $\mcT_{t} \circ \mcS$ is a variety, and $\{t(x,y) = x\}^p$ is its equational base.
\end{theorem}

\begin{corollary}\label{C:freePVt}
  For any subvariety $\mcV_t$ of $\mcT_t$, the free $\pro{\mcV_t}$-algebra
  $B = X\pro{\mcV_t}$ is a semilattice of $\mcV_t$-algebras. In particular,
  $B = \bigsqcup_{s \in S} B_s $, where $S = X\mcS$ is the free semilattice over $X$ and
  all $B_s$ are members of $\mcV_t$. Hence
  \begin{equation*}
B = X\pro{\mcV_t} \in \mcV_t \circ \mcS.
\end{equation*}
\end{corollary}

\begin{proposition}\label{P:quotientC}
  Let $\delta$ be a congruence of the free $\pro{\mcV_{t}}$-algebra $B = X\pro{\mcV_t}$
  over~$X$. Then the quotient $C = B/\delta$ is a semilattice sum of $\mcV_t$-algebras.
\end{proposition}

\begin{proof}
We first use the fact that $C$ is the quotient $B/\delta$ of $B$ to describe the
decomposition of $C$ into a semilattice sum of $\mcT_t$-algebras in terms of
the algebra $B$.
Let $\varrho_B$ be the semilattice replica congruence of the algebra $B$, and let
$\psi = \delta \vee \varrho_B$. Since $\psi \geq \varrho_B$ it follows that $B/\psi$ is a
semilattice. On the other hand, the minimality of $\varrho_B$ guarantees that $\psi$ is
the smallest congruence above $\delta$ such that $B/\psi \in \mcS$. It follows that
$\psi/\delta$ is the semilattice replica congruence of $B/\delta$.
Since each element of $C$ is of the form $b/\delta$ for some $b\in B$,
Theorem~\ref{T:varPt} implies that
\begin{equation*}
  C = \bigsqcup_{b/\delta\in B/\delta} (b/\delta)/(\psi/\delta)
\end{equation*}
in which each summand $(b/\delta)/(\psi/\delta) \in \mcT_t$.


On the other hand, by Corollary~\ref{C:freePVt}, $B = \bigsqcup_{s \in S} B_s $, in which
$S$ is a semilattice and all $B_s$ are members of $\mcV_t$.

It remains to show that the $(\psi/\delta)$-delta classes, $(b/\delta)/(\psi/\delta)$,
satisfy the identities of $\Gamma_t$.  Let
\begin{equation}\label{E:gamma}
x_1 \ldots x_k u = x_1 \ldots x_k v
\end{equation}
be a typical member of $\Gamma_t$.
We want to show that \eqref{E:gamma} holds in each $\psi/\delta$-class. This means that
for any $c^1 = b^1/\delta, \ldots, c^k = b^k/\delta \in (b/\delta)/(\psi/\delta)$, with
$b^i \in B_{s_i}$ and $b \in B$, we have
\begin{equation}\label{E:ginC}
c^1 \ldots c^k u = c^1 \ldots c^k v,
\end{equation}
which is equivalent to
\begin{equation}\label{E:ginB}
(b^1 \ldots b^k u, b^1 \ldots b^k v) \, \in \, \delta.
\end{equation}
Note that $b^i/\delta, b^j/\delta \in (b/\delta)/(\psi/\delta)$ if and only if
$(b^i, b^j) \in \psi$. Since each $(b^i /\delta)/(\psi/\delta)$ satisfies the identity
$x \cdot y = x$, we may use Lemma~\ref{L:equivlz} repeatedly to obtain the following:
\begin{multline*}
b^{i}\ \delta\ b^{i}\cdot b^{1}\ \delta \ (b^{i}\cdot b^{1})\cdot b^{2}\ \delta\ ((b^{i}\cdot b^{1})\cdot b^{2})\cdot b^{3}\ \delta\dots\\
\delta\ (\ldots((b^{i}\cdot b^{1})\cdot b^{2})\cdot \ldots)\cdot b^{k}\ =: d^i.
\end{multline*}
Thus $b^{i}\ \delta\ d^i$ for each $1 \leq i \leq k$, and the elements
$d^1, d^2, \ldots, d^k$ all belong to the same congruence class $B_{s_{1}\cdots s_k}$ of
$\varrho_B$.  Since $\varrho_B$-classes satisfy the identities of $\Gamma_t$, it follows
that
\begin{equation*}
d^1 \ldots d^k u = d^1 \ldots d^k v.
\end{equation*}
Finally note that
\begin{equation*}
(b^1 \ldots b^k u, d^1 \ldots d^k u) \in \delta \quad \mbox{and} \quad (b^1 \ldots b^k v,
d^1 \ldots d^k v) \in \delta,
\end{equation*}
which implies
\begin{equation*}
(b^1 \ldots b^k u, b^1 \ldots b^k v) \in \delta.
\end{equation*}
This verifies \eqref{E:ginB}.
\end{proof}

\end{document}
